\section{System}
\label{sec:system}

Figure~\ref{fig:pipeline} shows the pipeline. A perception agent classifies each heartbeat and emits an event. The events of a window cross an interface to a reasoning agent, which returns a structured triage decision. Only the interface varies between arms, and each arm implements one of the efficiency levers of Section~\ref{sec:intro} (Table~\ref{tab:levers}), so the comparison is one of levers on a fixed sender and a fixed receiver.

\begin{table}[t]
\caption{The interfaces compared, and the efficiency lever each implements. Prefix caching is applied to every arm as a second condition and is not a separate arm.}
\label{tab:levers}
\centering
\small
\begin{tabular}{lll}
\toprule
Arm & Lever & What crosses the interface \\
\midrule
A-full & none (reference) & Full record of every event, as text \\
A-compact & none (baseline) & One compact line per event, as text \\
A-filtered & say less & Abnormal events only, normal events as a count \\
Adapter (r4) & send latent & Four virtual tokens per event \\
\bottomrule
\end{tabular}
\end{table}

\begin{figure}[t]
\centering
\adjustbox{max width=\linewidth}{% Figure 1: system overview, included from sections/system.tex.
% Palette matches the data figures: orange = text with every event, violet = filtered text, blue = latent adapter.
% Grey fill = frozen component; thick coloured border + "TRAINED" badge = the one trained part.
\definecolor{cA}{HTML}{EB6834}
\definecolor{cF}{HTML}{4A3AA7}
\definecolor{cL}{HTML}{2A78D6}
\definecolor{cX}{HTML}{C0392B}
{\renewcommand{\baselinestretch}{1}\selectfont
\begin{tikzpicture}[font=\normalsize,
  frozen/.style={draw=black!55, rounded corners=3pt, align=center, inner sep=6pt, fill=black!6, text width=#1},
  lane/.style={draw=#1, line width=1pt, rounded corners=4pt, fill=#1!6, inner sep=0pt},
  tok/.style={draw=#1!80!black, fill=#1!35, minimum width=2.1mm, minimum height=4.2mm, inner sep=0pt, line width=0.4pt},
  arr/.style={-{Stealth[length=2.4mm]}, line width=0.9pt, black!60},
  lab/.style={font=\small, text=black!75, align=center, inner sep=1.5pt},
  badge/.style={font=\scriptsize\bfseries, text=white, fill=cL, rounded corners=2pt, inner sep=2.5pt}]

% ---- ECG strip with beat labels ------------------------------------------------------------------------
\def\qrs#1#2{\draw[line width=1pt,#2] (#1-0.45,0)--(#1-0.2,0)--(#1-0.1,0.18)--(#1,1.1)--(#1+0.12,-0.4)--(#1+0.24,0)--(#1+0.55,0);}
\node[font=\small\bfseries, anchor=south west] at (-0.2,1.9) {1. Sense};
\draw[black!35, line width=0.5pt] (-0.2,0)--(5.4,0);
\qrs{0.5}{black!80} \qrs{1.6}{black!80} \qrs{2.7}{cX} \qrs{3.8}{black!80} \qrs{4.9}{cX}
\foreach \x/\t/\c in {0.5/N/black!70,1.6/N/black!70,2.7/V/cX,3.8/N/black!70,4.9/S/cX}{\node[font=\small\bfseries,text=\c] at (\x,-0.75) {\t};}
\node[lab, anchor=north] at (2.6,-1.25) {$N$ consecutive heartbeats\\(here $N=5$)};

% ---- perception agent and events -------------------------------------------------------------------------
\node[font=\small\bfseries, anchor=south west] at (6.4,1.9) {2. Classify (frozen)};
\node[frozen=38mm] (per) at (8.3,0) {\textbf{Perception agent}\\CNN--LSTM + RR branch\\\emph{not a transformer}};
\draw[arr] (5.55,0)--(per.west);

\node[font=\small\bfseries, anchor=south west] at (11.9,1.9) {3. Report $N$ events};
\node[frozen=40mm] (ev) at (14.0,0) {context vector $\mathbf{z}_i\in\mathbb{R}^{32}$\\[1pt]label, confidence,\\heart rate, RR, run length};
\draw[arr] (per.east)--(ev.west);

% ---- three interfaces -------------------------------------------------------------------------------------
\node[font=\small\bfseries, anchor=south west] at (17.9,5.6) {4. Cross one of three interfaces};

% lane A: text, every event
\node[lane=cA, minimum width=74mm, minimum height=21mm] (la) at (22.0,3.6) {};
\node[font=\small\bfseries, text=cA!70!black, anchor=north west] at (18.3,4.6) {A-compact (baseline)};
\foreach \i in {0,...,17}{\node[tok=cA] at (18.65+0.27*\i,3.55) {};}
\node[lab, anchor=west] at (23.55,3.55) {$\ldots$};
\node[lab, anchor=north west] at (18.3,3.05) {every event as text: $\approx$60 tokens per event};

% lane B: filtered text
\node[lane=cF, minimum width=74mm, minimum height=21mm] (lb) at (22.0,0) {};
\node[font=\small\bfseries, text=cF!80!black, anchor=north west] at (18.3,1.0) {A-filtered (lever: say less)};
\foreach \i in {0,...,5}{\node[tok=cF] at (18.65+0.27*\i,-0.05) {};}
\node[lab, anchor=west, text=cX] at (20.35,-0.05) {V};
\foreach \i in {0,...,5}{\node[tok=cF] at (20.9+0.27*\i,-0.05) {};}
\node[lab, anchor=west, text=cX] at (22.65,-0.05) {S};
\node[draw=cF!80!black, fill=white, font=\scriptsize, inner sep=2pt, rounded corners=1pt] at (24.1,-0.05) {3 normal};
\node[lab, anchor=north west] at (18.3,-0.55) {abnormal events as text, normals as a count};

% lane C: latent adapter
\node[lane=cL, line width=2pt, minimum width=74mm, minimum height=21mm] (lc) at (22.0,-3.6) {};
\node[font=\small\bfseries, text=cL!80!black, anchor=north west] at (18.3,-2.6) {Adapter (lever: send latent)};
\node[badge, anchor=north east] at (25.5,-2.58) {TRAINED};
\foreach \g in {0,...,4}{\foreach \i in {0,...,3}{\node[tok=cL] at (18.65+1.3*\g+0.26*\i,-3.65) {};}}
\node[lab, anchor=north west] at (18.3,-4.15) {every event as 4 virtual tokens: $\mathbf{W}[\mathbf{z}_i;\mathbf{s}_i]+\mathbf{P}_i$};

\draw[arr] (ev.east)--++(0.9,0) |- (la.west);
\draw[arr] (ev.east)--(lb.west |- ev.east);
\draw[arr] (ev.east)--++(0.9,0) |- (lc.west);

% ---- reasoning agent -------------------------------------------------------------------------------------
\node[font=\small\bfseries, anchor=south west] at (26.9,5.6) {5. Decide (frozen)};
\node[frozen=36mm, minimum height=118mm] (llm) at (29.3,0) {\textbf{Reasoning agent}\\Gemma 4 E4B\\frozen, 4-bit\\[6pt]scaffold +\\schema-constrained\\decoding\\[8pt]$\downarrow$\\[2pt]\texttt{\{"tier": "urgent",}\\\texttt{ "beat": 3\}}};
\draw[arr] (la.east)--(llm.west |- la.east);
\draw[arr] (lb.east)--(llm.west |- lb.east);
\draw[arr] (lc.east)--(llm.west |- lc.east);

% ---- legend ----------------------------------------------------------------------------------------------
\node[frozen=0mm, minimum width=5mm, minimum height=4mm, inner sep=0pt] at (18.3,-6.2) {};
\node[lab, anchor=west] at (18.75,-6.2) {frozen (never updated)};
\node[draw=cL, line width=2pt, minimum width=5mm, minimum height=4mm, inner sep=0pt, fill=cL!6] at (22.4,-6.2) {};
\node[lab, anchor=west] at (22.85,-6.2) {the only trained part};
\node[draw=black!60, fill=black!25, minimum width=2.1mm, minimum height=4.2mm, inner sep=0pt, line width=0.4pt] at (26.4,-6.2) {};
\node[lab, anchor=west] at (26.75,-6.2) {= one token};
\end{tikzpicture}}
}
\caption{System overview. A perception agent that is not a transformer classifies each heartbeat, and the events of a window (abnormal beats in red) cross one of three interfaces to a frozen reasoning agent. Each small square is one token: text grows with every event (about 60 tokens each), the filter lists only abnormal events and counts the rest, and the adapter sends four virtual tokens per event. Grey boxes are frozen; only the adapter (blue) is trained. Sender, receiver and decoding are identical across arms, and the message sizes are the measured medians of Section~\ref{sec:results}.}
\label{fig:pipeline}
\end{figure}

\subsection{Task}
Each window contains $N$ consecutive heartbeats from one recording, with $N \in \{1,5,10,20,50\}$. The reasoning agent must return the most urgent triage tier in the window and say which beat set it. The tier of a single beat follows a fixed escalation rule applied to the perception agent's event: \emph{urgent} if the event requires urgent review (a run of three or more consecutive abnormal beats, or a single ventricular or fusion beat classified with confidence of at least 0.85), \emph{routine} if the beat is classified as normal, and \emph{priority} otherwise. The reference tier of a window is the maximum over its beats. The rule is applied to the perception agent's \emph{predictions}, not to the annotations, because the test is whether the channel delivers the sender's output faithfully; agreement with the annotations is reported separately (Section~\ref{sec:results}). The task is a communication-fidelity benchmark, not a clinical reasoning task. A fixed rule over the perception outputs solves it, so it cannot show that a language model adds clinical judgement.

\subsection{Perception Agent}
The perception agent is a convolutional--recurrent classifier trained only on the five-class AAMI EC57 beat task \cite{aami1998ec57}. A 360-sample window (1~s at 360~Hz) passes through three one-dimensional convolutional blocks (32, 64 and 128 channels, with batch normalisation, ReLU and max-pooling), a bidirectional LSTM with hidden size 64, and a unidirectional context LSTM with hidden size 32. The output of the context LSTM, a 32-dimensional \emph{context vector}, feeds the classification head and is the only representation passed downstream. There is no pretraining, no contrastive objective and no auxiliary loss.

The original design saw a single window and no timing, yet the text interface carried the RR interval, so a five-feature RR branch (previous and next RR interval, local mean RR and two ratios) was added to the context stage. Across ten training seeds it improved overall accuracy from 0.854 to $0.926 \pm 0.019$, ventricular sensitivity from 77\% to $93\% \pm 3\%$ and the class recoverability of the context vector from 69.6\% to $81.2\% \pm 5.0\%$, at the cost of one beat of latency. Fusion-beat detection stayed poor ($5\% \pm 7\%$ sensitivity). All adapters were trained on seed~0 of this encoder, which ranked third, fourth and seventh of ten on accuracy, supraventricular sensitivity and recoverability and was designated representative by a rule fixed beforehand.

A second, purely convolutional sender (residual one-dimensional convolutions without recurrence, the same RR branch and the same 32-dimensional vector) tests whether the channel depends on the first sender's architecture (Section~\ref{sec:results}).

\subsection{Reasoning Agent}
The reasoning agent is Gemma~4 E4B, loaded in 4-bit NF4 quantisation and frozen in every arm; neither adapter training nor calibration updates its weights. The model uses per-layer embeddings, a lookup keyed by token identity, which virtual tokens lack, so virtual positions receive the lookup of a padding token. This is an architectural cost of the latent channel and applies to every adapter result.

\subsection{Adapter: the Latent Channel}
The adapter is the third lever. Its single-event form, one linear layer from the 32-dimensional vector to four tokens with 337,920 trainable parameters, was described in the first author's MSc dissertation \cite{tasbir2026dissertation}. Here it is extended to many events and retrained, with the position embeddings, normalisation and side inputs described below. The setting imposes three requirements: the sender is unchanged and is not a transformer, the receiver is unchanged, and the message length depends only on the number of events. A linear map from the sender's own vector meets all three and keeps each event's tokens a function of that event alone. For event $i$ with context vector $\mathbf{z}_i \in \mathbb{R}^{32}$ and three scaled side inputs $\mathbf{s}_i \in \mathbb{R}^{3}$ (heart rate, RR interval and abnormal-run length, each clipped and standardised), the adapter forms $\mathbf{x}_i = [\mathbf{z}_i; \mathbf{s}_i] \in \mathbb{R}^{35}$ and produces $k=4$ virtual tokens
\begin{equation}
\mathbf{T}_i = \gamma \cdot \mathrm{norm}\!\left(\mathbf{W}\mathbf{x}_i + \mathbf{b} + \mathbf{P}_i\right) \in \mathbb{R}^{4 \times 2560},
\end{equation}
where $\mathbf{W}$ and $\mathbf{b}$ define one linear map from 35 to $4 \times 2560$ dimensions, $\mathbf{P}_i \in \mathbb{R}^{4\times 2560}$ is a learned embedding for event position $i$ (up to 50 positions), $\mathrm{norm}$ rescales each token to unit length, and $\gamma$ is one learnable scalar initialised to the median norm of the embeddings the model computes for real tokens. Without the normalisation the norm of the virtual tokens drifted far from that of real tokens, and many of them together broke generation. The window's tokens are the $\mathbf{T}_i$ concatenated in event order, in place of the events' text, so prompt length depends only on $N$.

Because the map is linear and acts on one event at a time, slot $i$ is a function of event $i$ alone, so each event's content can be audited independently (Section~\ref{sec:results}). The map has rank at most 35, so the adapter cannot introduce information absent from its input. The adapter has 880,641 trainable parameters (358,400 in the projection weights, 10,240 in its bias, 512,000 in the position embeddings and one scale), against the roughly four billion effective parameters of the frozen model.

Each window's training target is a JSON answer in the interface schema (the urgency tier, a templated justification and a guideline fact), built deterministically from the reference tier. Training uses teacher forcing: the correct answer's tokens follow the virtual tokens, and the loss is token-level cross-entropy over the answer only, $L = 4\,L_{\mathrm{tier}} + L_{\mathrm{rest}}$, where $L_{\mathrm{tier}}$ averages over the tokens that spell the tier and $L_{\mathrm{rest}}$ over the other target tokens. Gradients pass through the frozen language model to the adapter, and only the adapter's parameters are updated. We used AdamW at a learning rate of $5\times10^{-4}$, warm-up over the first 5\% of steps and cosine decay to 10\% of the peak, gradient clipping at 1.0, a batch of one window and at most three epochs, with validation every quarter epoch, early stopping from the second epoch, and the checkpoint with the best validation balanced accuracy. Recipe r4 trains on 447 windows spread over $N$ from 1 to 50 and validates on 170 windows from three held-out recordings. Figure~\ref{fig:training} shows the flow.

\begin{figure}[t]
\centering
\adjustbox{max width=\linewidth}{% Figure 2: how the adapter is trained, included from sections/system.tex.
% Grey fill = frozen; blue with thick border = trained; red = loss and the backward path (dashed).
\definecolor{cL}{HTML}{2A78D6}
\definecolor{cX}{HTML}{C0392B}
{\renewcommand{\baselinestretch}{1}\selectfont
\begin{tikzpicture}[font=\normalsize,
  frozen/.style={draw=black!55, rounded corners=3pt, align=center, inner sep=6pt, fill=black!6, text width=#1},
  io/.style={draw=black!40, rounded corners=3pt, align=center, inner sep=6pt, fill=white, text width=#1},
  train/.style={draw=cL, line width=2pt, rounded corners=3pt, align=center, inner sep=6pt, fill=cL!8, text width=#1},
  loss/.style={draw=cX, line width=1pt, rounded corners=3pt, align=center, inner sep=6pt, fill=cX!7, text width=#1},
  arr/.style={-{Stealth[length=2.4mm]}, line width=0.9pt, black!60},
  back/.style={-{Stealth[length=2.4mm]}, line width=1.1pt, dashed, cX!80},
  tok/.style={draw=cL!80!black, fill=cL!35, minimum width=2.1mm, minimum height=4.2mm, inner sep=0pt, line width=0.4pt},
  lab/.style={font=\small, text=black!75, align=center, inner sep=1.5pt}]

\node[io=40mm] (win) at (0,0) {\textbf{Training window}\\$N$ events, $N$ from 1 to 50,\\incl.\ hard-negative routine windows};
\node[frozen=40mm] (per) at (6.8,0) {\textbf{Perception agent}\\frozen\\gives $\mathbf{z}_i\in\mathbb{R}^{32}$ and side inputs $\mathbf{s}_i$};
\draw[arr] (win)--(per);

\node[train=60mm] (ad) at (14.6,0) {\textbf{Adapter}\\$\mathbf{T}_i=\gamma\,\mathrm{norm}(\mathbf{W}[\mathbf{z}_i;\mathbf{s}_i]+\mathbf{b}+\mathbf{P}_i)$\\880{,}641 trainable parameters};
\draw[arr] (per)--(ad);

\node[frozen=46mm, minimum height=26mm] (llm) at (24.0,0) {\textbf{Reasoning agent (frozen)}\\input: scaffold text, virtual tokens,\\then the correct answer tokens\\\emph{(teacher forcing)}};
\draw[arr] (ad.east)--(llm.west);
\foreach \i in {0,...,7}{\node[tok] at (18.85+0.27*\i,0.55) {};}
\node[lab, anchor=north] at (19.7,-0.05) {$4N$ virtual\\tokens};

\node[io=46mm] (tgt) at (24.0,5.4) {\textbf{Target answer}\\JSON: tier, templated justification,\\guideline fact\\(built from the reference tier)};
\draw[arr] (tgt)--(llm);

\node[loss=46mm] (ls) at (24.0,-5.0) {\textbf{Loss}\\$L = 4\,L_{\mathrm{tier}} + L_{\mathrm{rest}}$\\cross-entropy over the answer tokens only};
\draw[arr] (llm)--(ls);

\draw[back] (ls.west)--(ad.south |- ls.west)--(ad.south);
\node[lab, text=cX!80!black, anchor=south west] at (14.9,-4.6) {gradients pass through the frozen model;\\only the adapter is updated};

\node[io=44mm] (sel) at (6.8,-5.0) {\textbf{Schedule and selection}\\AdamW, lr $5\times10^{-4}$, batch of one;\\validation every quarter epoch,\\best balanced accuracy kept};
\end{tikzpicture}}
}
\caption{How the adapter is trained. The frozen perception agent supplies each event's vector and side inputs, the adapter (blue, the only trained part) turns them into virtual tokens, and the frozen reasoning agent is run on the scaffold, the virtual tokens and the correct answer (teacher forcing). The loss weights the tier tokens four times as heavily as the rest. Gradients (dashed) flow back through the frozen model but update only the adapter.}
\label{fig:training}
\end{figure} The final recipe, r4, differs from the earlier recipe r3 in two ways, both pre-registered before training. First, 72 training and 23 validation windows were added in which a normal beat is classified with confidence below 0.8 (\emph{hard negatives}), because r3 raised false alarms on exactly such windows. Second, the three side inputs were added, because heart rate, interval and run length were present in the text interface but could not be recovered from the 32-dimensional vector. The two changes were made together, and an ablation separates their effects (Section~\ref{sec:results}).

\subsection{Text Interfaces}
All text arms use the same scaffold, instructions and decoder as the adapter arm. They differ only in how the events are written into the prompt. A-compact and A-full list every event and so involve no effort to say less; A-filtered is the sender-side filter, the first lever.
\begin{itemize}
\item \textbf{A-compact:} one line per event with class, confidence, run length and signal quality. This is the strongest text arm that lists every event and is the primary comparator.
\item \textbf{A-full:} the full JSON record of each event, as in the earlier single-event study, reported only for cost.
\item \textbf{A-filtered:} only the non-normal events, with their position in the window, and the number of normal events as a count. No task information is lost: the filtered window yields the reference tier on all 1,195 test windows.
\end{itemize}
A text arm that lists every event is also reported under a \emph{calibrated} threshold (Section~\ref{sec:protocol}).

\subsection{Schema-Constrained Decoding}
Free generation produced malformed JSON in 9 to 20\% of answers across arms, through misnamed fields and loops. Every arm therefore uses the same constrained decoder. The structure of the answer is forced, the tier is restricted to the three valid words and chosen by the model's own scores, and free-text fields cannot close or break the structure. Every arm then produces a valid answer on every window. A check of 380 windows confirmed that the decoder changes no tier.
